\documentclass[11pt]{article}

\usepackage[margin=1in]{geometry}
\usepackage{amsmath,amssymb,booktabs,tabularx}

\title{News mark MGF}
\author{}
\date{}

\begin{document}
\maketitle

\noindent\textbf{Status.} Draft\\
\textbf{Plan item.} News mark MGF \(M_c(q,r)\)\\
\textbf{Depends on.} \texttt{03\_conditional\_cf\_mgf\_trading\_block.tex}\\
\textbf{Feeds.} Conditional CF/MGF of a trading block, affine Riccati system, affine domain constraints

\section{Purpose}

Derive the component-level transform for a coupled news mark.
The result is the mark MGF \(M_c(q,r)\) used when a news arrival
simultaneously contributes an activity shock and a return jump.

\section{Setup}

Fix a news component \(c\in\{\mathrm{known},\mathrm{surprise}\}\).
At one arrival time, draw a shared severity mark and a return-shock mark
\[
    U\sim\Gamma(k_U,\theta_U),
    \qquad
    \theta_U=\frac{1}{k_U},
    \qquad
    Z^N\sim\mathcal N(0,1),
\]
with \(U\) independent of \(Z^N\).
The gamma law uses the shape-scale parameterization,
so \(\mathbb E[U]=k_U\theta_U=1\).

The component scales are \(a_c\ge0\), \(m_c\in\mathbb R\), and
\(\sigma_c\ge0\).
The activity shock and return jump are
\[
    J^c=a_c U,
    \qquad
    \zeta^c=m_c U+\sigma_c\sqrt{U}\,Z^N.
\]
The transform arguments are \(r\), the loading on the activity shock, and
\(q\), the loading on the return jump.

\section{Notation}

\begin{tabularx}{\textwidth}{lXX}
\toprule
Symbol & Meaning & Units or domain\\
\midrule
\(c\) & News component & \(\{\mathrm{known},\mathrm{surprise}\}\)\\
\(U\) & Shared severity mark & \(\Gamma(k_U,1/k_U)\), non-negative\\
\(Z^N\) & Return-shock mark & Standard normal\\
\(a_c\) & Activity-coupling scale & Non-negative\\
\(m_c\) & Return-jump drift per severity unit & Real\\
\(\sigma_c\) & Return-jump volatility per square-root severity unit & Non-negative\\
\(J^c\) & Activity shock mark & Non-negative\\
\(\zeta^c\) & Coupled return-jump mark & Real\\
\(r\) & Activity-shock transform loading & Complex\\
\(q\) & Return-jump transform loading & Complex\\
\(M_c(q,r)\) & Joint mark MGF & Finite on the stated domain\\
\bottomrule
\end{tabularx}

\section{Assumptions used}

\begin{tabularx}{\textwidth}{XXX}
\toprule
Assumption & Where it is used & Why it is needed\\
\midrule
\(U\sim\Gamma(k_U,1/k_U)\) with \(k_U>0\).
& Evaluating the outer gamma MGF.
& Gives unit-mean severity and a closed-form mark transform.\\
\(Z^N\sim\mathcal N(0,1)\) and is independent of \(U\).
& Conditioning on \(U\) before evaluating the normal MGF.
& Produces the variance term \(\frac12 q^2\sigma_c^2 U\).\\
The same \(U\) drives \(J^c\) and \(\zeta^c\).
& Writing both marks inside one expectation.
& Encodes the shared-severity coupling between activity and returns.\\
The transform argument lies inside the gamma moment domain.
& Applying the gamma MGF.
& Ensures \(M_c(q,r)\) is finite.\\
\bottomrule
\end{tabularx}

\section{Result}

\noindent\textbf{Identity.} For a news component \(c\),
\[
    M_c(q,r)
    =
    \mathbb E\left[\exp\left(rJ^c+q\zeta^c\right)\right]
    =
    \left(
        1-\theta_U\left(ra_c+qm_c+\frac{1}{2}q^2\sigma_c^2\right)
    \right)^{-k_U}.
\]

Equivalently, the log mark MGF is
\[
    \ell_c(q,r)
    =
    \log M_c(q,r)
    =
    -k_U\log\left(
        1-\theta_U\left(ra_c+qm_c+\frac{1}{2}q^2\sigma_c^2\right)
    \right).
\]

\section{Derivation}

Start from the joint transform of the activity shock and return jump:
\[
M_c(q,r)
=
\mathbb E\left[
    \exp\left(ra_cU+q\left(m_cU+\sigma_c\sqrt{U}\,Z^N\right)\right)
\right].
\]
Condition on \(U\):
\[
    M_c(q,r)
    =
    \mathbb E\left[
        \exp\left((ra_c+qm_c)U\right)
        \mathbb E\left[
            \exp\left(q\sigma_c\sqrt{U}\,Z^N\right)\mid U
        \right]
    \right].
\]
For a standard normal \(Z^N\),
\[
    \mathbb E\left[
        \exp\left(q\sigma_c\sqrt{U}\,Z^N\right)\mid U
    \right]
    =
    \exp\left(\frac{1}{2}q^2\sigma_c^2U\right),
\]
where the formula is interpreted by analytic continuation for complex
arguments within the transform domain.

Therefore
\[
    M_c(q,r)
    =
    \mathbb E\left[
        \exp\left(
            \left(ra_c+qm_c+\frac{1}{2}q^2\sigma_c^2\right)U
        \right)
    \right].
\]
Let
\[
    s_c(q,r)
    =
    ra_c+qm_c+\frac{1}{2}q^2\sigma_c^2.
\]
The gamma MGF for \(U\sim\Gamma(k_U,\theta_U)\) is
\[
    \mathbb E[e^{sU}]=(1-\theta_U s)^{-k_U},
    \qquad
    \operatorname{Re}(s)<\frac{1}{\theta_U}.
\]
Substituting \(s=s_c(q,r)\) gives
\[
    M_c(q,r)
    =
    \left(1-\theta_U s_c(q,r)\right)^{-k_U}
    =
    \left(
        1-\theta_U\left(ra_c+qm_c+\frac{1}{2}q^2\sigma_c^2\right)
    \right)^{-k_U}.
\]

\section{Domain and validity}

The result requires \(k_U>0\), \(\theta_U=1/k_U\), \(a_c\ge0\),
\(\sigma_c\ge0\), and finite \(m_c\).
The transform is valid when
\[
    \operatorname{Re}\left(
        ra_c+qm_c+\frac{1}{2}q^2\sigma_c^2
    \right)
    <
    \frac{1}{\theta_U}.
\]
For the affine schedule transform, \(r\) is supplied by the Hawkes loading,
typically \(r=\mathbf D(\tau)\cdot\mathbf r\), and this domain condition must
hold along the Riccati path for every active news component.

\section{Consequences}

\begin{itemize}
    \item The same closed-form \(M_c(q,r)\) handles activity-only, return-only, and joint news-mark queries.
    \item The shared severity creates transform coupling through the single scalar \(s_c(q,r)\).
    \item A component with \(a_c=0\) and \(m_c=\sigma_c=0\) contributes \(M_c(q,r)=1\), so inactive news marks do not alter the transform.
    \item The affine domain constraint for news marks is exactly the gamma MGF moment-domain constraint evaluated along the schedule.
\end{itemize}

\section{Checks}

\begin{tabularx}{\textwidth}{XXX}
\toprule
Check & Expected result & Evidence location\\
\midrule
Gamma severity normalization
& Sampled severity marks have mean one and second moment \(1+1/k_U\).
& Test \texttt{gamma\_severity\_distribution\_regression\_matches\_unit\_mean\_moments} in \texttt{crates/janiform-model-kernels/tests/news\_coupling\_distribution\_regression.rs}\\
Gamma-normal construction
& The constructed activity shock is \(a_cU\), and the constructed return jump is \(m_cU+\sigma_c\sqrt U Z^N\).
& Test \texttt{gamma\_normal\_news\_mark\_couples\_activity\_and\_return\_to\_shared\_severity} in \texttt{crates/janiform-model-kernels/tests/news\_coupling\_distribution\_regression.rs}\\
Component validation
& Activity coupling is non-negative, return volatility is non-negative, and return jumps without activity coupling are rejected.
& \texttt{NewsComponentMarkScaleConfiguration::validate} in \texttt{crates/janiform-model-kernels/src/news/configuration.rs}\\
Affine use
& The schedule transform uses \(M_c(q,\mathbf D(\tau)\cdot\mathbf r)\) and \(\ell_c(q,\mathbf D(\tau)\cdot\mathbf r)\).
& \texttt{design/janiform\_compact\_model\_spec.tex}\\
\bottomrule
\end{tabularx}

\end{document}
